\section{Related Work}
\label{sec:related}

\subsection{LLM-Powered Browser Agents}

The field of LLM-based web agents has grown rapidly. Browser-Use~\cite{browseruse} provides a Playwright-based agent framework that integrates vision-language models for web interaction. SeeAct~\cite{seeact} pioneered the use of GPT-4V for integrated visual understanding and action grounding, achieving 51.1\% on live websites with manual grounding. WebRL~\cite{webrl} introduced self-evolving online curriculum RL, improving Llama-3.1-8B from 4.8\% to 42.4\% on WebArena-Lite. CUA-Gym~\cite{cuagym} demonstrated scalable RL with verifiable rewards, reaching 72.6\% on OSWorld-Verified. OpenSkynet differs from these systems in its focus on \emph{post-execution improvement}: rather than improving single-task performance through training, it improves cross-task robustness through self-healing and skill reuse at inference time.

\subsection{Skill Reuse and Memory}

Several recent works address skill reuse and memory in agent systems. MUSE-Autoskill~\cite{museautoskill} introduces a skill-centric framework with a unified skill lifecycle (creation, memory, management, evaluation, refinement). MEMENTO~\cite{memento} treats the web as a learning signal, with dual-channel memory (declarative knowledge + procedural strategies). ExpeL~\cite{expel} extracts insights from successful and failed trajectories via LLM reflection. OpenSkynet's skill learning pipeline is inspired by these approaches but adds iterative refinement with failure-aware context and similar-skill merging. The three-tier memory model with heuristic importance scoring and background consolidation is novel in the browser automation space.

\subsection{Self-Healing and Error Recovery}

Self-healing has been explored in GUI automation (RoSUE~\cite{rosue} repairs UI test scripts) and robotic task execution (Robo-Cortex~\cite{robocortex} maintains a Navigation Heuristic Library). CRITIC~\cite{critic} validates agent outputs with external tools before committing. OpenSkynet's healer extends these ideas to browser automation by combining multimodal evidence (screenshot + DOM) with LLM reasoning to produce targeted step-level patches, versioning each fix for auditability.

\subsection{Agent Decomposition and Planning}

Plan-and-execute architectures~\cite{react} decompose tasks into subtasks for focused execution. ReAct~\cite{react} interleaves reasoning and action in a single loop. Reflexion~\cite{reflexion} adds verbal reinforcement learning through stored reflections. MiRA~\cite{mira} uses subgoal-driven planning with milestone-based rewards. OpenSkynet's Manager Agent combines task classification, strategy selection, and beam-width decomposition; the subagent system extends these ideas with permission-based isolation and controlled parallel execution.

\subsection{Benchmarks for Browser Agents}

WebArena~\cite{webarena} and VisualWebArena provide reproducible environments for evaluating web agents. Mind2Web~\cite{mind2web} tests generalization to unseen websites. WebVoyager~\cite{webvoyager} evaluates on live websites. The Evil Skill Generalization Benchmark proposed in this work is complementary: it focuses on \emph{adversarial robustness} and \emph{skill generalization under parameter variation}, dimensions not systematically tested by existing benchmarks.
